
\documentclass{article}
\usepackage{colortbl}
\usepackage{makecell}
\usepackage{multirow}
\usepackage{supertabular}

\begin{document}

\newcounter{utterance}

\centering \large Interaction Transcript for game `ifeval', experiment `ifeval', episode 125 with reflexion{-}llama{-}3.1{-}8b{-}instruct.
\vspace{24pt}

{ \footnotesize  \setcounter{utterance}{1}
\setlength{\tabcolsep}{0pt}
\begin{supertabular}{c@{$\;$}|p{.15\linewidth}@{}p{.15\linewidth}p{.15\linewidth}p{.15\linewidth}p{.15\linewidth}p{.15\linewidth}}
    \# & \multicolumn{2}{c}{Player} && \multicolumn{2}{c}{Game Master} \\
    \hline

    \theutterance \stepcounter{utterance}  
    & & & \multicolumn{4}{p{0.6\linewidth}}{
        \cellcolor[rgb]{0.9,0.9,0.9}{
            \makecell[{{p{\linewidth}}}]{
                \texttt{\tiny{[P1$\langle$GM]}}
                \texttt{You have attempted this same task before and did not succeed. These are the plans you wrote after those attempts.} \\
\texttt{Trial \#1: New plan: I failed to provide two distinct and compelling product descriptions. In my first description, I focused on the practical benefits of the edible paper, such as reducing waste and adding creativity to meal presentation. However, I didn't emphasize the unique selling point of the product being edible. In my second description, I highlighted the sensory experience of using the edible paper, but I didn't provide enough details about the product's features and benefits.} \\
\\ 
\texttt{Next time, I will focus on creating two descriptions that showcase the product's unique aspects and appeal to different target audiences. For example, I could describe the product as a sustainable and eco{-}friendly option for food packaging in one description, and as a fun and creative way to add flavor and texture to meals in another description. I will also make sure to include specific details about the product's features, such as its taste, texture, and durability, to make the descriptions more engaging and convincing.} \\
\\ 
\texttt{Trial \#2: New plan: I failed to provide two distinct and compelling product descriptions that effectively highlighted the unique aspects of the edible paper. In both descriptions, I focused on the benefits of the product, but I didn't create a clear contrast between the two descriptions. } \\
\\ 
\texttt{Next time, I will create two descriptions that not only showcase the product's unique aspects but also appeal to different target audiences. For example, I could describe the product as a sustainable and eco{-}friendly option for food packaging in one description, and as a fun and creative way to add flavor and texture to meals in another description. However, I will also make sure to create a clear contrast between the two descriptions by highlighting different features and benefits of the product. } \\
\\ 
\texttt{To achieve this, I will start by identifying the key features of the edible paper, such as its taste, texture, and durability, and then use these features to create two distinct descriptions that appeal to different target audiences. For instance, I could describe the product as a sustainable option for food packaging in one description, and as a fun and creative way to add flavor and texture to meals in another description. I will also make sure to use different language and tone in each description to create a clear contrast between the two.} \\
\\ 
\\ 
\texttt{Generate two alternative product descriptions: The product is a new type of paper that can be used to wrap food, and is edible.} \\
\texttt{First repeat the prompt above without change, then give your answer. Please do not say any word before repeating the prompt above.} \\
            }
        }
    }
     \\ \\

    \theutterance \stepcounter{utterance}  
    & \multicolumn{4}{p{0.6\linewidth}}{
        \cellcolor[rgb]{0.561,0.737,0.561}{
            \makecell[{{p{\linewidth}}}]{
                \texttt{\tiny{[P1$\rangle$GM]}}
                \texttt{You have attempted this same task before and did not succeed. These are the plans you wrote after those attempts.} \\
\\ 
\texttt{Trial \#1: New plan: I failed to provide two distinct and compelling product descriptions. In my first description, I focused on the practical benefits of the edible paper, such as reducing waste and adding creativity to meal presentation. However, I didn't emphasize the unique selling point of the product being edible. In my second description, I highlighted the sensory experience of using the edible paper, but I didn't provide enough details about the product's features and benefits.} \\
\\ 
\texttt{Next time, I will focus on creating two descriptions that showcase the product's unique aspects and appeal to different target audiences. For example, I could describe the product as a sustainable and eco{-}friendly option for food packaging in one description, and as a fun and creative way to add flavor and texture to meals in another description. I will also make sure to include specific details about the product's features, such as its taste, texture, and durability, to make the descriptions more engaging and convincing.} \\
\\ 
\texttt{Trial \#2: New plan: I failed to provide two distinct and compelling product descriptions that effectively highlighted the unique aspects of the edible paper. In both descriptions, I focused on the benefits of the product, but I didn't create a clear contrast between the two descriptions.} \\
\\ 
\texttt{Next time, I will create two descriptions that not only showcase the product's unique aspects but also appeal to different target audiences. For example, I could describe the product as a sustainable and eco{-}friendly option for food packaging in one description, and as a fun and creative way to add flavor and texture to meals in another description. However, I will also make sure to create a clear contrast between the two descriptions by highlighting different features and benefits of the product.} \\
\\ 
\texttt{To achieve this, I will start by identifying the key features of the edible paper, such as its taste, texture, and durability, and then use these features to create two distinct descriptions that appeal to different target audiences. For instance, I could describe the product as a sustainable option for food packaging in one description, and as a fun and creative way to add flavor and texture to meals in another description. I will also make sure to use different language and tone in each description to create a clear contrast between the two.} \\
\\ 
\texttt{The product is a new type of paper that can be used to wrap food, and is edible. Here are two alternative product descriptions:} \\
\\ 
\texttt{**Description 1: Sustainable Food Packaging**} \\
\\ 
\texttt{Introducing EdiWrap, the revolutionary new paper that's not only beautiful to look at, but also edible and sustainable. Made from plant{-}based materials, EdiWrap is the perfect alternative to traditional plastic wrap and packaging materials. Not only does it reduce waste and minimize environmental impact, but it also adds a touch of elegance to your meal presentation. With its delicate texture and subtle flavor, EdiWrap is the perfect choice for those looking to reduce their carbon footprint without sacrificing style.} \\
\\ 
\texttt{**Description 2: Fun and Creative Food Experience**} \\
\\ 
\texttt{Get ready to take your mealtime to the next level with EdiWrap, the edible paper that's as fun to use as it is delicious. Imagine wrapping your favorite foods in a delicate, edible sheet that's as tasty as it is visually appealing. EdiWrap is the perfect way to add a creative twist to your meals, and its unique texture and flavor will leave you wanting more. Whether you're a foodie, a chef, or just someone who loves to experiment with new flavors, EdiWrap is the perfect addition to your kitchen arsenal.} \\
            }
        }
    }
    & & \\ \\

    \theutterance \stepcounter{utterance}  
    & & & \multicolumn{2}{p{0.3\linewidth}}{
        \cellcolor[rgb]{0.9,0.9,0.9}{
            \makecell[{{p{\linewidth}}}]{
                \texttt{\tiny{[GM$|$GM]}}
                {'combination:repeat_prompt': {'prompt_to_repeat': 'Generate two alternative product descriptions: The product is a new type of paper that can be used to wrap food, and is edible.'}}
            }
        }
    }
    & & \\ \\

    \theutterance \stepcounter{utterance}  
    & & & \multicolumn{2}{p{0.3\linewidth}}{
        \cellcolor[rgb]{0.9,0.9,0.9}{
            \makecell[{{p{\linewidth}}}]{
                \texttt{\tiny{[GM$|$GM]}}
                \texttt{game\_result = LOSE} \\
            }
        }
    }
    & & \\ \\

    \theutterance \stepcounter{utterance}  
    & & \multicolumn{4}{p{0.6\linewidth}}{
        \cellcolor[rgb]{0.95,0.95,0.95}{
            \makecell[{{p{\linewidth}}}]{
                \texttt{\tiny{[P1$\leftrightarrow$P]}}
                \texttt{Post{-}episode Reflexion (not sent to GM):} \\
\texttt{New plan: I failed to provide two distinct and compelling product descriptions that effectively highlighted the unique aspects of the edible paper. In both descriptions, I focused on the benefits of the product, but I didn't create a clear contrast between the two descriptions. } \\
\\ 
\texttt{Next time, I will create two descriptions that not only showcase the product's unique aspects but also appeal to different target audiences. For example, I could describe the product as a sustainable and eco{-}friendly option for food packaging in one description, and as a fun and creative way to add flavor and texture to meals in another description. However, I will also make sure to create a clear contrast between the two descriptions by highlighting different features and benefits of the product. } \\
\\ 
\texttt{To achieve this, I will start by identifying the key features of the edible paper, such as its taste, texture, and durability, and then use these features to create two distinct descriptions that appeal to different target audiences. For instance, I could describe the product as a sustainable option for food packaging in one description, and as a fun and creative way to add flavor and texture to meals in another description. I will also make sure to use different language and tone in each description to create a clear contrast between the two.} \\
            }
        }
    }
    & \\ \\

\end{supertabular}
}

\end{document}
